\noindentparagraph{\bf{Limitations and Future Work.}}
In RPWS, time multiplexing is used to realize mathematically exact alpha blending, hence many frames are required for accurate reconstruction. In practice, the effectiveness of time multiplexing is constrained by SLM refresh rates, though this limitation may be alleviated by advances in high-speed SLMs, such as the 3600\,fps FLCoS used in Holographic Parallax~\cite{kim2024holographic} and the recently developed 5760\,fps TI MEMS-based PLM. Addressing quantization artifacts in these devices presents an interesting future direction. Our experiments also show reduced contrast due to random phase, but extending learning-based calibration techniques~\cite{peng2020neural, choi2022time} to directly reconstruct random-phase complex fields~\cite{jang2024waveguide} may further enhance RPWS holograms. Please refer to the supplemental materials for more discussions on contrast enhancement strategies.


\noindentparagraph{\bf{Conclusion.}} Random-phase Wave Splatting enables photorealistic CGH with correct defocus blur and parallax, unlocking perceptually realistic VR. It generalizes beyond Gaussians to diverse translucent primitives, seamlessly bridging modern neural scene representations with next-generation display technology.
